\subsection{局部凸可度量化空间的对偶}

命题 2. --- 设 E 是一个局部凸可度量化空间，F 是它的强对偶。空间 F 是完备的、半桶型的，且满足下列条件：

(DB) 存在有界子集的一个序列 \((\mathbf{A}_n)_{n\in \mathbb{N}}\)（\(\mathbf{F}\) 中的），使得 \(\mathbf{F}\) 的每个有界子集都包含在某个 \(\mathbf{A}_n\) 中。

空间 E 是有界型的（III，第 12 页，命题 2），从而它的强对偶是完备的（III，第 23 页，推论 1）。

设 \((\mathrm{V}_{n})_{n\in\mathbf{N}}\) 是 E 中零邻域的一个递减序列，使得 E 的每个零邻域都包含某个 \(V_{n}\)。设 \(A_{n}\) 是 \(V_{n}\) 在 F 中的极集。由于 E 是有界型的，F 的每个有界子集都等度连续（III，第 22 页，命题 10），从而包含在某个 \(A_{n}\) 中。换言之，空间 F 满足条件 (DB)。

我们现在证明 F 是半桶型的。设 \((\mathrm{U}_{n})_{n\in\mathbf{N}}\) 是 F 中凸、均衡且闭的零邻域的一个序列。我们假定集 \(U=\bigcap_{n}U_{n}\) 吸收 F 的每个有界子集。我们将证明 U 是 F 中 0 的一个邻域。为此，我们将对整数 \(n\geqslant0\) 用归纳法构造实数 \(\lambda_{n}>0\) 与 F 中凸、均衡的零邻域 \(W_{n}\)（它们对于 \(\sigma(\mathrm{F},\mathrm{E})\) 是闭的），使之满足下列关系

\[
\lambda_ {n} \mathrm{A} _ {n} \subset \frac {1}{2} \mathrm{U} \cap \left(\bigcap_ {0 \leqslant i <   n} \mathrm{W} _ {i}\right)\tag{1}
\]

\[
\bigcup_ {0 \leqslant i \leqslant n} \lambda_ {i} \mathrm{A} _ {i} \subset \mathrm{W} _ {n} \subset \mathrm{U} _ {n}.\tag{2}
\]

假定诸数 \(\lambda_{i}\) 与诸集 \(\mathrm{W}_i\) 已对 \(0 \leqslant i < n\) 构造出来。根据假设，集 U 吸收 F 的有界子集；此外，对 \(0 \leqslant i < n\)，\(\mathrm{W}_i\) 是 F 中 0 的一个邻域，从而吸收 F 的有界子集。因此可以找到满足 (1) 的数 \(\lambda_{n} > 0\)。设 C 表示对于 \(\sigma(\mathrm{F}, \mathrm{E})\) 而言 \(\bigcup_{0 \leqslant i \leqslant n} \lambda_i \mathrm{A}_i\) 的闭凸均衡包；集 C 是等度连续的，

从而对于 \(\sigma(\mathrm{F},\mathrm{E})\) 紧（III，第 17 页，推论 2）。由于 \(\mathrm{U}_n\) 是 F 中 0 的一个邻域，存在 E 的有界子集 B 使得 \(\mathbf{B}^{\circ}\subset\frac{1}{2}\mathbf{U}_{n}\)。令 \(\mathrm{W}_n=\mathrm{C}+\mathrm{B}^{\circ}\)。由于 \(\mathbf{B}^{\circ}\) 是 F 中 0 的一个邻域，我们看出 \(\mathrm{W}_n\) 是 F 中凸、均衡的零邻域。此外，C 紧且 \(\mathbf{B}^{\circ}\) 对于 \(\sigma(\mathrm{F},\mathrm{E})\) 闭；由 GT, III, § 4, 第 1 目的推论 1，\(\mathrm{W}_n\) 对于 \(\sigma(\mathrm{F},\mathrm{E})\) 闭。最后，我们有 \(\mathbf{C}\subset\frac{1}{2}\mathbf{U}\subset\frac{1}{2}\mathbf{U}_n\) 与 \(\mathbf{B}^{\circ}\subset\frac{1}{2}\mathbf{U}_n\)，从而 \(\mathrm{W}_n\subset\mathrm{U}_n\)，因为 \(\mathrm{U}_n\) 是凸的。这样我们就建立了 (2)。

令
\[
\mathrm{W} = \bigcap_ {n} \mathrm{W} _ {n}
\]

则
\[
\mathbf {W} \subset \mathbf {U}
\]

由 (1) 与 (2)，我们有
\[
\lambda_ {i} \mathrm{A} _ {i} \subset \mathrm{W} _ {j} \quad \text {for all } i \text { and } j
\]

在 N 中，从而 \(\lambda_{i}A_{i}\subset W\)（对所有 \(i\in N\)）。特别地，W 是吸收的，从而是对于 \(\sigma(\mathrm{F},\mathrm{E})\) 的一个桶。由 IV 第 4 页的注 3，W 是 F 中 0 的一个邻域。更加地，U 是 F 中 0 的一个邻域，故 F 是半桶型的。

下列推论把 Banach-Steinhaus 定理推广到 Fréchet 空间的对偶（参见 III，第 25 页，推论 2）。

推论. --- 设 G 是一个 Hausdorff 局部凸空间，\((u_{n})\) 是从 F 到 G 中的线性映射的一个序列，它简单收敛于一个从 F 到 G 中的映射 u。则 u 连续，且序列 \((u_{n})\) 在 F 的每个预紧子集上一致收敛于 u。

由于 F 是完备的，所有 \(u_{n}\) 组成的集（它对于简单收敛拓扑有界）在 \(\mathcal{L}_{b}(F;G)\) 中有界（III，第 27 页，推论 1）。由于空间 F 是半桶型的（命题 2），\(\mathcal{L}_{b}(F;G)\) 的每个可数且有界的子集由 IV 第 21 页的命题 1 都是等度连续的。因此诸 \(u_{n}\) 组成的集等度连续，且推论由 III 第 18 页的推论得出。

